% Source: ACM_Templates.pdf, PDF pp. 82-83; use standard Euler-tour RMQ.
\subsection{欧拉序与 ST 表：常数时间 LCA}
\label{tree:rmq-lca}
\AlgorithmInfo
    {对固定树进行大量 LCA 查询，单次查询为常数时间。}
    {ST 表预处理 $O(n\log n)$，查询 $O(1)$，空间 $O(n\log n)$。}
    {非空连通树，顶点从 1 开始；只比较深度，边权不参与计算。}
\subsubsection*{原理}
进入一个点时记录它，每次从儿子返回再记录当前点，得到长度 $2n-1$ 的欧拉序。记每个点首次出现的位置为 $first_u$，则 $u,v$ 的 LCA 是欧拉序上两次首次出现位置之间\textbf{深度最小的点}。因为 DFS 必须经过 LCA 才能从一侧走到另一侧，而这段遍历不会越过更浅的祖先。

对欧拉序建立按深度取最小点的 ST 表，长度为 $L$ 的区间可由两个长 $2^{\lfloor\log_2 L\rfloor}$ 的区间覆盖。最小值允许重叠，合并两个结果即可在 $O(1)$ 内回答。常数时间指查询，不是这份 ST 实现的预处理复杂度。
\subsubsection*{使用教程}
\InlineCode{build(root)} 后用 \InlineCode{lca(u,v)} 查询。\InlineCode{ord} 是从 0 开始的欧拉序，\InlineCode{first[u]} 是其下标；原树顶点仍从 1 开始。新增边或换根后重新构建。
\begin{lstlisting}[style=mtycpp,numbers=none]
RMQLCA g;
g.init(4);
g.add(1,2),g.add(1,3),g.add(2,4);
g.build(1);
int p=g.lca(3,4);
\end{lstlisting}
\needspace{7\baselineskip}
\subsubsection*{模板代码}
\CodeFile{code/trees/rmq_lca.cpp}
